\documentclass[10pt,letterpaper]{article}
\usepackage{technical-testing}
\setDocTitle{Ballistics And Spring Problem Set}
\setDocSubtitle{Original problems with numerical checks and evidence requirements}
\setDocVersion{1.0}
\setDocStatus{Educational / Draft}
\setDocOwner{Game Development}
\setDocDate{2026-10-30}
\setDocSummary{Original problems with numerical checks and evidence requirements}
\begin{document}
\makedoctitle
\makedocsummary

\section{Problems}
\begin{enumerate}
\item Given $y_0=0$, $v_0=6$ m/s and $g=-9.81$ m/s$^2$, compute the time to apex and maximum height. Sample the curve every $1/120$ s and compare with the analytical result.
\item Compare three spring damping ratios: 0.2, 1.0, and 2.0, holding natural frequency constant. Which settles fastest without overshoot in an appropriately stable simulation?
\item Show why the same force produces a different acceleration when mass doubles. Then retime two silhouette animations to convey different weights without changing drawn mass.
\item Demonstrate a large-step Euler integration failure and justify the chosen fixed step.
\end{enumerate}
\section{Reference check}
The analytical apex occurs at $t=6/9.81\approx0.612$ s and $y_{\max}=6^2/(2\cdot9.81)\approx1.835$ m. The visual motion may exaggerate this result, but the physics check should identify that departure explicitly.
\section{Evidence}
Store the model and its coordinate conventions, sample file, discrepancy analysis, graph, and an intentionally stylized variant.
\end{document}
